\chapter{Regularization and Ensemble Methods}
\label{ch:regularization}

\epigraph{``The best single model is rarely as good as a collection of mediocre models.''}{}

\learninggoal{
	\item Understand regularization as a mechanism to control model complexity and prevent overfitting.
	\item Derive ridge regression and lasso, and contrast L1 vs.\ L2 regularization.
	\item Explain how bagging reduces variance without increasing bias.
	\item Understand boosting as adaptive reweighting of misclassified points.
	\item Build random forests and gradient-boosted trees for structured data.
}

\section{Regularization}

Regularization adds a penalty to the loss function to discourage overly complex models:

\[
	L_{\text{reg}}(\mathbf{w}) = L(\mathbf{w}) + \lambda \cdot P(\mathbf{w})
\]

\subsection{Ridge Regression (L2)}

\[
	P(\mathbf{w}) = \|\mathbf{w}\|_2^2 = \sum_{j=1}^d w_j^2
\]

\[
	\hat{\mathbf{w}}_{\text{ridge}} = (\mathbf{X}^\top \mathbf{X} + \lambda \mathbf{I})^{-1} \mathbf{X}^\top \mathbf{y}
\]

Ridge shrinks coefficients toward zero but never to exactly zero. It works well when many features have small effects.

\begin{definition}[Bias--Variance in Ridge]
	Ridge increases bias (coefficients are biased toward zero) but decreases variance (coefficients are more stable). The total prediction error can decrease if the variance reduction outweighs the bias increase~\cite{hoerl1970ridge}.
\end{definition}

\subsection{Lasso (L1)}

\[
	P(\mathbf{w}) = \|\mathbf{w}\|_1 = \sum_{j=1}^d |w_j|
\]

Lasso drives some coefficients to exactly zero, performing automatic feature selection~\cite{tibshirani1996regression}. This makes it useful when only a few features are relevant.

\subsection{Elastic Net}

Combines L1 and L2 penalties:

\[
	P(\mathbf{w}) = \alpha \|\mathbf{w}\|_1 + (1 - \alpha) \|\mathbf{w}\|_2^2
\]

Elastic net handles groups of correlated features better than lasso (which tends to pick one from a group).

\subsection{The Regularization Path}

As $\lambda$ increases from 0 to $\infty$:

\begin{itemize}
	\item Ridge: coefficients smoothly approach 0.
	\item Lasso: coefficients shrink linearly; some hit 0 at finite $\lambda$.
\end{itemize}

Cross-validation selects $\lambda$ to minimize validation error.

\section{Bagging}

Bootstrap Aggregating (bagging) reduces variance by averaging many models trained on bootstrap samples:

\begin{enumerate}
	\item Draw $B$ bootstrap samples from the training data.
	\item Train a model on each sample.
	\item Average predictions (regression) or majority vote (classification).
\end{enumerate}

\begin{theorem}[Variance Reduction]
	If $B$ models each have variance $\sigma^2$ and pairwise correlation $\rho$, the ensemble variance is:
	\[
		\Var(\bar{f}) = \rho \sigma^2 + \frac{1 - \rho}{B} \sigma^2
	\]
	As $B \to \infty$, variance converges to $\rho \sigma^2$. Bagging is most effective when models are uncorrelated ($\rho \approx 0$).
\end{theorem}

\section{Random Forests}

Random forests improve on bagged trees by decorrelating the trees~\cite{breiman2001random}:

\begin{itemize}
	\item Each tree is trained on a bootstrap sample.
	\item At each split, only a random subset of $m$ features is considered ($m \approx \sqrt{d}$ for classification, $d/3$ for regression).
	\item This decorrelates the trees, reducing $\rho$ and thus variance.
\end{itemize}

\begin{lstlisting}[language=Python, caption=Random forest]
from sklearn.ensemble import RandomForestRegressor

rf = RandomForestRegressor(
    n_estimators=500,
    max_features='sqrt',
    min_samples_leaf=5,
    n_jobs=-1
)
rf.fit(X_train, y_train)
\end{lstlisting}

Random forests are among the best off-the-shelf methods for structured (tabular) data. They handle non-linearity, interactions, missing values, and mixed data types.

\section{Boosting}

Boosting trains models sequentially, each focusing on the mistakes of the previous ones.

\subsection{AdaBoost}

Adaptive Boosting assigns weights to training points:
\begin{enumerate}
	\item Initialize weights $w_i = 1/n$.
	\item For $t = 1, \dots, T$:
	      \begin{itemize}
		      \item Train a weak learner $f_t$ on weighted data.
		      \item Compute weighted error $\epsilon_t$.
		      \item Set $\alpha_t = \frac12 \log\frac{1 - \epsilon_t}{\epsilon_t}$.
		      \item Increase weights for misclassified points: $w_i \gets w_i e^{\alpha_t \mathbbm{1}[y_i \neq f_t(\mathbf{x}_i)]}$.
	      \end{itemize}
	\item Ensemble: $\hat{y} = \text{sign}(\sum_t \alpha_t f_t(\mathbf{x}))$.
\end{enumerate}

\subsection{Gradient Boosting}

Generalizes boosting to any differentiable loss function:

\[
	f_{t+1}(\mathbf{x}) = f_t(\mathbf{x}) + \eta \cdot h_t(\mathbf{x})
\]

where $h_t$ is a weak learner (typically a shallow tree) fit to the \emph{negative gradient} of the loss:

\[
	h_t \approx -\nabla_{f} L(y, f_t(\mathbf{x}))
\]

For squared error loss, the gradient is just the residual $y - f_t(\mathbf{x})$, so gradient boosting is ``fitting to the errors''.

\subsection{XGBoost, LightGBM, CatBoost}

Modern gradient boosting implementations add:

\begin{itemize}
	\item \textbf{XGBoost~\cite{chen2016xgboost}:} regularization, second-order gradients, weighted quantile sketch.
	\item \textbf{LightGBM:} Gradient-based One-Side Sampling (GOSS), Exclusive Feature Bundling (EFB), leaf-wise tree growth.
	\item \textbf{CatBoost:} ordered boosting (handles categorical features), symmetric trees.
\end{itemize}

\begin{lstlisting}[language=Python, caption=XGBoost]
import xgboost as xgb

model = xgb.XGBRegressor(
    n_estimators=1000,
    learning_rate=0.01,
    max_depth=6,
    subsample=0.8,
    colsample_bytree=0.8
)
model.fit(X_train, y_train,
          eval_set=[(X_val, y_val)],
          early_stopping_rounds=50)
\end{lstlisting}

\section{Stacking}

Stacked generalization combines multiple different models via a meta-learner:

\begin{enumerate}
	\item Split training data into $K$ folds.
	\item For each base model $m$, generate out-of-fold predictions for all training points.
	\item Train a meta-model on these predictions as features.
\end{enumerate}

Stacking can improve over any single model but risks overfitting if not done carefully.

\section{Practical Guidelines}

\begin{longtable}{p{3cm}p{5cm}p{5cm}}
	\toprule
	\textbf{Scenario}                   & \textbf{Recommended method} & \textbf{Why}               \\
	\midrule
	$n \gg d$, linear relationship      & Ridge or Lasso              & Simple, interpretable      \\
	$d \gg n$, sparse signal            & Lasso                       & Feature selection          \\
	Tabular data, $n \sim 10^3$--$10^5$ & XGBoost / LightGBM          & Best predictive accuracy   \\
	Heterogeneous data types            & Random forest               & Robust, no tuning needed   \\
	Need calibrated probabilities       & Gradient boosting + Platt   & Better calibration than RF \\
	Very large $n$ ($>10^6$)            & Linear model or LightGBM    & Computational constraints  \\
	\bottomrule
\end{longtable}

\begin{principle}[Swets' Law]
	The performance of a well-tuned tree ensemble and a well-tuned neural network on tabular data is often similar. The main difference is in computational cost, interpretability, and data efficiency.
\end{principle}

\section{Exercises}

\begin{exercise}
	Fit ridge regression on a dataset with correlated features. Plot the coefficient paths as $\lambda$ varies. Compare with lasso: which coefficients go to zero first? Why?
\end{exercise}

\begin{exercise}
	Train a single decision tree, a random forest with 100 trees, and an XGBoost model on the same dataset. Compare training error and test error. Which ensemble reduces variance most? Which reduces bias most?
\end{exercise}

\begin{exercise}
	Implement AdaBoost from scratch (or using scikit-learn's \texttt{AdaBoostClassifier}) with decision stumps (depth-1 trees) as weak learners. Plot the training and test error as a function of the number of boosting rounds.
\end{exercise}

\begin{exercise}
	Explain why random forests use both bootstrap sampling and random feature subsets. What would happen if you used only one of these sources of randomness?
\end{exercise}
